%% =====================================================================
%%  PKNU-HCMUT Joint International Academic Symposium on Engineering
%%  Mini paper  --  main file
%%
%%  Topic   : Design and Implementation of a Real-Time Flight Controller
%%            Architecture on ESP32-S3 for Quadcopter UAVs
%%  Derived from: 261030_PKNU-HCMUT-Abstract_Real-Time Flight Controller.docx
%%  Content plan: ../poster_draft_v2.md  (abstract rewritten per its v3 rules)
%%
%%  All layout lives in style.sty; each section body lives in
%%  01..06_*.tex (one file per chapter) and is pulled in with \input.
%%  Paragraphs are kept on single logical lines (LaTeX treats a newline as
%%  a space; soft wrap in the editor handles the display).
%%
%%  Build   : .\build.ps1   (pdflatex x2)
%%            or Zed task "pdflatex: Build Paper (2 Passes)"
%% =====================================================================
\documentclass[10pt,a4paper]{article}

% --- Custom Style Package --------------------------------------------
\usepackage{style}

% --- Document Information --------------------------------------------
\symposium{PKNU-HCMUT Joint International Academic Symposium on Engineering}
\title{Design and Implementation of a Real-Time Flight Controller Architecture on ESP32-S3 for Quadcopter UAVs}
\author{Quy Hao Nhien Nguyen\textsuperscript{1,2}, Gia-Huy Luu\textsuperscript{1,2}, Cat Gia Han Nguyen\textsuperscript{1,2}, Ngo Dai Cat Truong\textsuperscript{1,2}, Van-Thanh Nguyen\textsuperscript{1,2,*}, Quoc-Nguyen Banh\textsuperscript{1,2}}
\affiliation{\textsuperscript{1}Faculty of Mechanical Engineering, Ho Chi Minh City University of Technology (HCMUT), Ho Chi Minh City, Vietnam\\ \textsuperscript{2}Vietnam National University Ho Chi Minh City, Ho Chi Minh City 70000, Vietnam}
\corresponding{nvthanh@hcmut.edu.vn}

\begin{document}

\maketitle

% =====================================================================
%  Abstract  (kept in the main file: it is part of the title block)
%  Rewritten from poster_draft_v2.md: architecture-as-contribution framing,
%  honesty protocol, only §R-cleared numbers + [NUM] slots for Step 0.
% =====================================================================
\paperheading{Abstract}

\begin{leadblock}
Community projects have demonstrated that ESP32-class system-on-chips can fly quadcopters; the open gap is not feasibility but evidence. In PX4 the ESP32 appears only as a WiFi telemetry module, the ArduPilot ESP32 port remains experimental, Espressif's ESP-Drone has seen limited support since 2022, and no published platform couples a system-level determinism measurement with a verifiable architecture ruleset. This paper presents FC-V01, a five-layer flight-controller architecture for a 450~mm quadcopter on a dual-core 240~MHz ESP32-S3 under FreeRTOS, whose dependencies point strictly downward from application to platform HAL and whose invariants stay checkable through decision records and gated audits (155 cross-layer includes scanned; three known upward edges, all under closure). The firmware performs no dynamic allocation in flight, decouples multi-rate tasks with lock-free ring buffers, and accelerates IMU ellipsoid calibration through PIE SIMD primitives on a self-designed four-layer carrier with a cost-optimised sensor set, sustaining a 500~Hz inertial acquisition and estimation chain. Determinism is measured rather than assumed: SystemView timing and jitter statistics [NUM: full-system profile pending Step 0], sensor-noise characterisation, and logging-robustness probes are reported from bench validation with the vehicle tethered, including the altitude estimate reported as limited by the low-grade BMP180 barometer. Free flight has not yet been performed; an INDI control law is future work.
\end{leadblock}

\keywordline{ESP32-S3; FreeRTOS; real-time determinism; layered software architecture; Extended Kalman Filter; low-cost UAV.}

% =====================================================================
%  Sections (one file per section, in reading order)
% =====================================================================
% =====================================================================
%  1 Introduction
%  Source: poster_draft_v2.md ① (~80-100 words base, expand to ~350-400)
%  - Context: hard real-time chain DAQ -> preprocessing -> estimation -> control
%  - Gap (cite §R1-R4, R7): PX4 has no ESP32 board; ArduPilot ESP32 port is
%    "experimental"; ESP-Drone limited support since 12/2022.
%  - Claim: feasibility exists; the gap is measured determinism + architectural
%    rules -> the contribution of this work.
%  - Fold the §RW positioning into ONE short paragraph here (keep the
%    comparative table for Discussion).
% =====================================================================
\section{Introduction}

% =====================================================================
%  2 Methodology  (largest block — keep ~1 page per half)
%  Source: poster_draft_v2.md ②
% =====================================================================
\section{Methodology}

% ---------------------------------------------------------------------
%  2.1 Hardware
%  - ESP32-S3 dual-core @ 240 MHz (R14), BLE, native SIMD (PIE via compute HAL)
%  - GY-87: MPU6050 (accel+gyro), HMC5883L (mag), BMP180 (baro)
%  - MicroSD blackbox, WS2812B LEDs, LoRa footprint, power-monitor board
%  - 4-layer KiCad PCB FC_v01 (FC_v02 in progress); gerbers + BOM released (R15)
%  - Figure: PCB render/photo; end with the design-intent sentence
% ---------------------------------------------------------------------
\subsection{Hardware}

% ---------------------------------------------------------------------
%  2.2 Firmware  (centre of the paper)
%  - Five-layer general architecture L5 app -> L4 core -> L3 middleware ->
%    L2 drivers -> L1 platform HAL; dependencies point DOWNWARD only (strict
%    downward DAG); decision records + gated ATAM-style audits
%    (diagram: FC_v01-general-architecture-v2.html)
%  - RTOS task set / periods / core assignment / IPC
%    (diagram: FC_v01-rtos-architecture.html; 14 resident tasks at boot,
%    CALIB_* and BLE_OTA spawn on demand — R13)
%    -> present as Table 1 (task | period | core), data from config.h (R8)
%  - SIMD: hal_simd primitives (dot/add/sub/mul/scale/mat_mul/inv3x3) for
%    float, packed int16, packed int8 -> ESP-DSP (dsps_*/dspm_*) -> S3 PIE
%    SIMD; hal_matrix type-safe facade; zero dynamic allocation;
%    beneficiaries: EKF prediction/update, calibration fits, filtering
% ---------------------------------------------------------------------
\subsection{Firmware}

% ---------------------------------------------------------------------
%  2.3 Experimental Setup  (static analyses — how you measured)
%  - Sensor measurements: time-domain histogram/stdev per axis; FFT -> noise
%    floor / interference bands  [NUM: dated log]
%  - Task periods & jitter: nominal from config.h vs SystemView JTAG/TRAX
%    trace (NOT SD); median/p99/max per task, >=1e4 loops, full load + RF
%    [NUM: pending Step 0 re-profile]
%  - Logging robustness: BLE MTU validation (currently never checked), SD
%    fsync latency (3-30 ms) vs ring depth  [NUM: counters + histogram]
%  - Estimation: attitude EKF 7-state quaternion + gyro bias, dual-task async
%    500/100 Hz; altitude EKF 3-state (baro+accel) 100 Hz
%    [NUM: altitude 24.15% is OFFLINE REPLAY ONLY — say so explicitly]
%  - Calibration: gyro static offset; accel 6-DOF ellipsoid (6 orientations);
%    mag 9-DOF ellipsoid; SIMD-accelerated, zero dynamic RAM (keep short)
%  - End with the mandatory protocol note: bench validation, tethered/clamped;
%    free flight NOT yet performed; INDI is future work
% ---------------------------------------------------------------------
\subsection{Experimental Setup}

% =====================================================================
%  3 Results  (one number per item + source)
%  Source: poster_draft_v2.md ③
%  - System frequency table + core map from SystemView (replaces the
%    theoretical ~20-22% load claim; note-0007 data = early DAQ-only phase,
%    use only with stated scope — R10, R16)
%    Static table usable now: IMU DAQ 500 Hz / preprocess 500 Hz / EKF
%    roll-pitch 500 Hz / EKF yaw 100 Hz / EKF altitude 100 Hz / FSM 100 Hz /
%    baro 50 Hz / mag 100 Hz / RC 50 Hz / motor 250 Hz / SD+BLE 100 Hz;
%    DAQ budget = 2.0 ms
%  - IMU charts: stdev histogram + FFT per axis  [NUM: dated log]
%  - Attitude EKF capability  [NUM: accuracy vs reference + per-task timing]
%  - Altitude EKF — told truthfully: performance is POOR; cause = BMP180
%    (low-grade, old-gen I2C baro), presented as a designed-in hardware
%    limitation with the improvement path in Discussion
%    [NUM: 24.15% = OFFLINE REPLAY ONLY (note 0037)]
%  - SIMD calibration: residuals + fit time, zero dynamic RAM  [NUM]
%  - End-to-end logging: SD blackbox 100 Hz + BLE telemetry + BLE OTA
%  - Honesty box (mandatory): "Measured on bench; free flight not yet performed."
% =====================================================================
\section{Results}

% =====================================================================
%  4 Discussion
%  Source: poster_draft_v2.md ④
%  - Why the altitude EKF underperforms: BMP180 trades resolution / rate /
%    noise for cost; the estimator cannot recover information the sensor
%    never delivered. Improvement path: higher-resolution higher-rate digital
%    baro OR ToF + optical flow -> FC_v02 hardware work
%  - Architecture robustness -> flyability: five-layer downward-only DAG held
%    through two gated audits (155 cross-layer includes scanned; only 3 known
%    upward edges, all with closure plans; C3 cycle caught and closed — R9).
%    Determinism is measured, not assumed (SystemView + jitter over >=1e4
%    loops) -> the architecture is a verifiable invariant, not a drawing
%  - Position vs related work: no platform publishes a full-system determinism
%    trace; closest (arXiv 2607.23220) claims "jitter-free" without a profiler
%    trace. -> condense the §RW table here (arXiv 2607.23220, IEEE iSES 2021,
%    IJPEDS 2024, UFABC/CBA, esp-drone) with a "how we differ" column
% =====================================================================
\section{Discussion}

% =====================================================================
%  5 Conclusion and Contributions
%  Source: poster_draft_v2.md ⑤ (~150 words, 4 numbered contributions)
%  1. Five-layer RTOS architecture for an ESP32-S3 FC with MEASURED
%     determinism (SystemView), 500 Hz DAQ + EKF
%  2. Zero-dynamic-allocation design + zero-RAM SIMD calibration — technique
%     set for RAM-constrained MCUs
%  3. Evidence-driven process: decision records + gated architecture audits —
%     reusable model for embedded projects
%  4. Self-designed hardware: 4-layer carrier PCB (FC_v01 -> FC_v02) +
%     power-monitor board, gerbers and BOM released
%  Future work: INDI controller (slot 0034), mixer, baro upgrade / ToF +
%  optical flow (altitude-hold path), free flight, full-system + RF re-profiling.
%  Repeat the honesty statement ONCE here.
% =====================================================================
\section{Conclusion and Contributions}

% =====================================================================
%  References
%  Source: poster_draft_v2.md REFERENCES (tier-filtered) + §RW tier note
%  Primary: PX4 docs, ArduPilot docs, Espressif, arXiv, IEEE, IJPEDS, SBA/CBA.
%  Low-tier (IJRASET/IJRTI/IJSREM) only as existence proofs.
%  Plan: switch to a .bib file (e.g. references.bib + \bibliography) once the
%  list stabilises; until then keep \begin{thebibliography}{10} here.
%
%  1.  PX4 Dev Team, "ESP32 WiFi Module", docs.px4.io/main/en/telemetry/esp32_wifi_module
%  2.  PX4 Dev Team, "WiFi Telemetry Radio", docs.px4.io/main/en/telemetry/telemetry_wifi
%  3.  PX4 Dev Team, "Flight Controller (Autopilot) Hardware", docs.px4.io/main/en/flight_controller/
%  4.  ArduPilot Dev Team, "ESP32 Based Autopilot", ardupilot.org/dev/docs/esp32-autopilot.html
%  5.  Espressif, "ESP-Drone", github.com/espressif/esp-drone (GPL-3.0; Crazyflie-derived)
%  6.  Mahbub, Roy, Prince, "Embedded Firmware Development for Flight Control, Telemetry, and Video Streaming for DIY UAV Research", arXiv:2607.23220, 2026
%  7.  Parikh, Vohra, Kaveshgar, "Comparison of Attitude Estimation Algorithms With IMU Under External Acceleration", IEEE iSES 2021
%  8.  Moussaoui et al., "Experimental study of PID for attitude control of a quadcopter using an ESP32", IJPEDS 15(3), 2024
%  9.  da Silva, Fazzolari, Correa, "Design and Implementation of a Real Time Attitude Estimation System with Low Cost Sensors", SBA/CBA (UFABC)
%  10. FreeRTOS (Kernel + SMP), POSA vol.1 (Layers), Clean Architecture — methodology references
% =====================================================================
\section*{References}


\end{document}
